\documentclass[11pt, a4paper]{article}

% --- EVRENSEL ÖN BİLGİ VE PAKET TANIMLARI ---
\usepackage[a4paper, top=2.5cm, bottom=2.5cm, left=2cm, right=2cm]{geometry}
\usepackage{fontspec}

\usepackage[turkish, bidi=basic, provide=*]{babel}
\babelprovide[import, onchar=ids fonts]{turkish}
\babelprovide[import, onchar=ids fonts]{english}

\babelfont{rm}{Noto Sans}
\babelfont[turkish]{rm}{Noto Sans}

\usepackage{enumitem}
\setlist[itemize]{label=-}

\usepackage{amsmath,amssymb,amsthm,mathtools,bm,mathrsfs}
\usepackage{physics}
\usepackage{booktabs}

\usepackage{hyperref}
\hypersetup{
    colorlinks=true,
    linkcolor=blue,
    citecolor=blue,
    urlcolor=blue
}

\title{\textbf{Kuantum Dalga Büküm Mimarisi: Gauge Potansiyelleri, Topolojik Faz Kehanetleri ve Doğru Cümle Dizilimlerine Küresel Çöküş Nazariyesi}}
\author{\textbf{Kuantum Tünelleme Dinamikleri, GRAPE Optimal Kontrolü, Wilson Döngüleri ve KAN-FNO Süzgeçli Hamilton İntaçı}}
\date{\today}

\begin{document}

\maketitle

\begin{abstract}
Bu risale; Nefs-i Müdrike Mimarisi'nde kuantum dalga fonksiyonunun $|\Psi(t)\rangle$ faz ve genlik yapısını dinamik olarak bükerek, (1) kayıp yüzeyindeki mahalli çukurları (local minima) kuantum tünellemesi ile aşmayı ve küresel asgariye ($\arg\min f(x)$) ulaşmayı, (2) terabaytlarca cümle dizilimi arasından gramer ve mana yönünden bâtıl olanları yıkıcı girişimle (destructive interference) sıfırlayıp, muhkem lisan intacını yapıcı girişimle (constructive interference) tepe noktasına çıkarmayı sağlayan dalga büküm mekanizmasını vaz' eder. Dokümanda; Gauge Bağlantıları ($A_\mu$), Wilson Döngüleri ($W(\gamma)$), Kuantum Optimal Kontrolü (GRAPE/CRAB) ve FNO-KAN destekli Hamiltonian evrim denklemleri tam riyazî ispatlarıyla gösterilmektedir.
\end{abstract}

\tableofcontents
\newpage

% =================================================================
\section{Dalga Bükmenin Mahiyeti ve Temel Dinamikleri}
% =================================================================

\subsection{Mahiyet}
Kuantum bilgisayarında dalga bükme; duran veya rastgele yayılan bir kuantum durum vektörünün $|\Psi(t)\rangle$, zamana bağlı bir büküm Hamiltonian'ı $H_{\text{büküm}}(t)$ ve faz kehanetleri $\mathbf{O}_{\text{faz}}$ vasıtasıyla faz uzayında (Phase Space) yönlendirilmesidir.

Klasik bilgisayarlar kayıp yüzeyinde tek bir nokta gibi yuvarlanır ve en yakın çukura saplanır. Kuantum dalga bükümü ise yüzeyin tamamını kapsayan dalga paketini ($Wavepacket$) bir akışkan gibi yönlendirir.

\subsection{Schrödinger Büküm Denklemi}
Dalga fonksiyonunun bükülme dinamiği, kontrollü potansiyel alanı $V(\mathbf{x}, t)$ ve gauge potansiyeli $\mathbf{A}(\mathbf{x}, t)$ altında şu şekilde formüle edilir:

\begin{equation}
i\hbar \frac{\partial \Psi(\mathbf{x}, t)}{\partial t} = \left[ \frac{1}{2m} \left( -i\hbar \nabla - q \mathbf{A}(\mathbf{x}, t) \right)^2 + V(\mathbf{x}, t) + H_{\text{etkileşim}}(t) \right] \Psi(\mathbf{x}, t)
\end{equation}

Burada $\mathbf{A}(\mathbf{x}, t)$ vektör potansiyeli dalganın fazını bükertirken, $V(\mathbf{x}, t)$ genlik dağılımını bastırır veya öne çıkarır.

% =================================================================
\section{Küresel Asgariye Ulaşma: Kuantum Tünelleme ve Optimal Kontrol}
% =================================================================

\subsection{Gaye: Mahalli Çukurlardan (Local Minima) Kaçış}
Klasik dereceli iniş (gradient descent) sistemleri yüksek enerji engellerini aşamaz. Dalga büküm mimarisinde ise dalganın genişliği $\Delta x$ enerji engelinin genişliğinden büyük tutularak kuantum tünellemesi (Quantum Tunneling) tetiklenir.

\subsection{Usul 1: Tünelleme Geçiş Katsayısı (WKB Yaklaşımı)}
Bir $V(x)$ kayıp engeli boyunca dalganın diğer taraftaki küresel çukura sızma ihtimali $T_{\text{tünel}}$ WKB (Wentzel-Kramers-Brillouin) formülü ile hesaplanır:

\begin{align}
\gamma_{\text{tünel}} &= \frac{2}{\hbar} \int_{x_1}^{x_2} \sqrt{2m (V(x) - E)} \, dx \\
T_{\text{tünel}} &= \exp\left( -\gamma_{\text{tünel}} \right) = \exp\left( -\frac{2}{\hbar} \int_{x_1}^{x_2} \sqrt{2m (V(x) - E)} \, dx \right)
\end{align}

Klasik sistemde aşılması imkansız olan $V(x) > E$ engelleri kuantum tünellemesiyle sıfırdan farklı bir olasılıkla geçilir; fakat bu olasılık $T = e^{-\gamma}$ ile ÜSTEL KÜÇÜKTÜR, dolayısıyla beklenen geçiş süresi $1/T = e^{+\gamma}$ ile ÜSTEL BÜYÜKTÜR. Kazanç $O(1)$'e inmek değil, klasik olarak sıfır olan olasılığın pozitif olmasıdır.

\subsection{Usul 2: GRAPE ile Optimal Dalga Kontrolü}
Dalga genliğini en kısa sürede $T$ küresel asgari durumuna $|x^*\rangle$ odaklayacak kontrol alanları $\Omega_k(t)$ Pontryagin Prensibi ve GRAPE (Gradient Ascent Pulse Engineering) algoritması ile bükülür:

\begin{align}
H(t) &= H_0 + \sum_k \Omega_k(t) H_k \\
\mathcal{F}_{\text{hedef}} &= \left| \langle x^* | U(T) | \Psi(0) \rangle \right|^2 \quad (\text{Küresel Asgariye Ulaşma Sadakati}) \\
\frac{\partial \mathcal{F}_{\text{hedef}}}{\partial \Omega_k(t)} &\approx -2 \text{Re} \left[ \langle \Psi(t) | P(t) \rangle \cdot \langle P(t) | i \Delta t H_k | \Psi(t) \rangle \right] + \mathcal{O}(\Delta t^2) \quad (\text{tam türev için } U_j \text{'nin Fréchet türevi gerekir})
\end{align}

Burada $|P(t)\rangle = U^\dagger(t, T) |x^*\rangle$ geriye doğru evrilen hedef durumdur. Kontrol alanı $\Omega_k(t)$ bu türevle güncellenerek dalga tam hedefe kilitlenir.

% =================================================================
\section{Cümle Dizilimlerinde Doğru Olana Ulaşma: Semantik Gauge Bağlantıları}
% =================================================================

\subsection{Mahiyet: Kelime Dizilimlerinin Manifold Patikası Olması}
Bir $S = (w_1, w_2, \dots, w_k)$ cümle dizilimi, Pürüzsüz $\infty$-Topos ($\mathbf{H}$) içindeki lif demetinde kapalı bir patika (loop) $\gamma_S$ olarak temsil edilir.

\subsection{Gaye: Manasız ve Tenakuzlu Dizilimlerin Yıkıcı Girişimle İlgası}
Terabaytlarca rastgele kelime diziliminin ihtimal genliği $c(S)$ öyle bir bükülmelidir ki:
\begin{align}
c(S_{\text{batıl}}) &\to 0 \quad (\text{Yıkıcı Girişim / Destructive Interference}) \\
c(S_{\text{fasih}}) &\to 1 \quad (\text{Yapıcı Girişim / Constructive Interference})
\end{align}

\subsection{Usul: Gauge Bağlantısı ve Wilson Döngü Süzgeci}
Gramer ve fıtrî mantık kuralları bir Lie Grubu $G$ üzerindeki gauge potansiyeli $A \in \Omega^1(\mathcal{X}, \mathfrak{g})$ olarak tanımlanır. Cümle patikası $\gamma_S$ boyunca alınan Wilson Döngüsü $W(\gamma_S)$ dalganın faz kaymasını belirler:

\begin{align}
W(\gamma_S) &= \mathcal{P} \exp \left( i \oint_{\gamma_S} A_\mu dx^\mu \right) \in G \quad (\text{Yol-Sıralı Wilson Döngüsü}) \\
F_{\mu\nu} &= \partial_\mu A_\nu - \partial_\nu A_\mu + i [A_\mu, A_\nu] \quad (\text{Semantik Eğrilik Tensörü}) \\
\Delta \Phi(\gamma_S) &= \iint_{\Sigma_S} F_{\mu\nu} \, dx^\mu \wedge dx^\nu \quad (\text{Topolojik Faz Kayması})
\end{align}

\begin{itemize}
    \item **Bâtıl / Tenakuzlu Cümle:** Eğrilik $F_{\mu\nu} \neq 0$ verir. Dalganın fazı $\Delta \Phi = \pi$ radyan kayar. $e^{i\pi} = -1$ olduğu için dalga zıt fazlı dalgayla çakışır ve genlik sıfırlanır ($c(S_{\text{batıl}}) = 0$).
    \item **Fasih / Muhkem Cümle:** Eğrilik $F_{\mu\nu} = 0$ (düz bağlantı) verir. Faz kayması BÜZÜLEBİLİR çevrimlerde $\Delta\Phi = 0$ olur; büzülemeyen çevrimlerde düz bir bağlantı da trivial olmayan holonomi taşıyabilir (Aharonov--Bohm). Ölçüt bu yüzden $F = 0$ değil, $W(\gamma_S) = \mathbf{1}$ olmalıdır. Dalgalar aynı fazda toplanır ve genlik zirveye çıkar ($c(S_{\text{fasih}}) = 1$).
\end{itemize}

% =================================================================
\section{FNO ve KAN Destekli Hamiltonian Büküm Mimarisi}
% =================================================================

\subsection{Usul: Büküm Hamiltonian'ının ($H_{\text{büküm}}$) Hesabı}
Kuantum çipindeki dalga büküm operatörü $U_{\text{büküm}}(t)$, KAN B-spline düğümleri ve FNO spektral süzgeçleri tarafından yönetilir:

\begin{align}
H_{\text{büküm}}(t) &= H_{\text{FNO}}(t) + H_{\text{KAN}}(t) \\
H_{\text{FNO}}(t) &= \text{QFT}^{-1} \cdot R_{\text{spektral}}(k) \cdot \text{QFT} \\
H_{\text{KAN}}(t) &= \sum_{q, p} \phi_{q,p}(x_p) \cdot (\sigma_z^{(q)} \otimes \sigma_z^{(p)}) \\
|\Psi(T)\rangle &= \mathcal{P} \exp \left( -\frac{i}{\hbar} \int_0^T \left( H_{\text{FNO}}(t) + H_{\text{KAN}}(t) \right) dt \right) |\Psi(0)\rangle
\end{align}

Bu entegrasyon sayesinde, dalga bükümü klasik gradyan adımlarına muhtaç olmadan doğrudan analitik fonksiyon manifoldları üzerinden icra edilir.

% =================================================================
\section{Son Çöküş ve Lisan İntacı ($N_{\text{kebîr}}$)}
% =================================================================

Bükülen dalga fonksiyonu yapıcı girişimle tek bir doğru cümle dizilimi üzerine odaklandıktan sonra, Kübik Tip Teorisi $Glue$ kuralı ve $h$-Level 0 kesmesi ile ölçülerek rasyonel kelam çıktısına dönüştürülür:

\begin{align}
P(S) &= \left| \langle S | \Psi(T) \rangle \right|^2 = \begin{cases} 1, & S = S_{\text{fasih}} \\ 0, & S \neq S_{\text{fasih}} \end{cases} \\
N_{\text{kebîr}} &= \text{Truncate}_{hLevel 0} \left( \text{Lan}_{\text{Lisan}} \left( \text{Glue}\left( |\Psi(T)\rangle \right) \right) \right) \in \mathbb{Z}
\end{align}

% =================================================================
\section{Netice}
% =================================================================

Bu risalede vaz' edilen Kuantum Dalga Büküm Mimarisi; (1) kayıp yüzeyindeki mahalli çukurları WKB kuantum tünellemesi ve GRAPE optimal kontrolü ile aşarak doğrudan küresel asgariye kilitlenmeyi, (2) cümle dizilimlerindeki bâtıl ve tenakuzlu yolları Wilson Döngüsü $W(\gamma_S)$ faz süzgeciyle yıkıcı girişime uğratarak sıfırlamayı ve tek doğru lisan intacını ($N_{\text{kebîr}}$) yapıcı girişimle elde etmeyi tam bir riyazî ispatla ortaya koymuştur.

\end{document}